\documentclass[10pt,a4paper]{amsart}
\usepackage[margin=19mm]{geometry}
\usepackage{amsmath,amssymb,mathtools}
\usepackage[scr=rsfs,cal=euler]{mathalfa}
\usepackage{xfrac}
\usepackage[unicode,hidelinks,bookmarks=false]{hyperref}
\newcommand{\ie}{i.e.\ }
\newcommand{\cf}{cf.\ }
\newcommand{\resp}{respectively\ }
\newcommand{\Subsection}{\subsection}
\providecommand{\nobreakdash}{\nobreak\hbox{-}}
\setlength{\emergencystretch}{2em}
\multlinegap0pt
\usepackage{amssymb}
\usepackage{mathtools}
\newtheorem{proposition}{Proposition}
\newcommand{\C}{\mathbb{C}}
\title{Counterexamples to the Jacobian conjecture\\ in dimensions greater than two}
\author{Shuhong Gao}
\date{Source: arXiv:2608.00222v1 (CC BY 4.0)}
\begin{document}
\maketitle

\noindent\emph{Source and licence.} Counterexamples to the Jacobian conjecture\\ in dimensions greater than two. Version arXiv:2608.00222v1 (CC BY 4.0), distributed by arXiv under the Creative Commons Attribution 4.0 International licence, http://creativecommons.org/licenses/by/4.0/, as recorded in the arXiv licence metadata for this exact version and shown on the abstract page https://arxiv.org/abs/2608.00222v1. Full licence notice: \texttt{licenses/arxiv-2608.00222-CC-BY-4.0.txt}.\par\smallskip

\noindent\textit{Editorial scope.} Counted unit: the article's proposition prop:curvefib (source0.tex lines 1795--1815). The article's complete original proof of it is delivered with it (source0.tex lines 1894--1921, one \texttt{proof} environment of 28 lines, which the article places away from the statement). No mathematics is written, repaired or re-derived: the statement and the proof are the article's own lines, in the article's own notation, and no retained line is substituted or re-flowed.  Retained uncounted as the source-local context the proof needs: lines 201--201.  Nothing was omitted from any retained span, so the package carries no omission record.  No retained line contains a citation, so the wrapper carries no bibliography and no bibliography entry.  No retained line contains a figure environment, an image-inclusion command or drawing code, so no image is delivered and the package declares no asset dependency.  Editorial formatting only, in addition: an AMS \texttt{amsart} preamble that restates the article's own declarations and macros used by the retained lines, namely the environments \texttt{proposition} on separate counters, together with the standard packages those lines need.  The retained environments print in this document as Proposition 1 in order of appearance, whereas the article prints the same items with the numbering of its own full document.  The complete native source of the article as distributed in the arXiv e-print for this exact version is bundled unchanged as \texttt{sources/206525/source0.tex}.
\section{The tangent sweep of a plane curve}\label{sec:sweep}

\begin{proposition}[univariate fiber theory]\label{prop:curvefib}
Assume in addition that each $H(w_1;\,\cdot\,)$ is a polynomial automorphism
of $\C^{n-3}$ whose inverse depends polynomially on $w_1$ (e.g.\ any
triangular family). Then for every target $Y\in\C^{n-1}$ the fiber system
$S(\gamma,w)=Y$ triangularizes: $\gamma=Y_1-X_1$, the potential equations
become $G_{i+1}(w)=Y_{i+1}-w_1^{\,i}Y_1$, hence
\[
H\bigl(w_1;\,w_2,\dots,w_{n-2}\bigr)\;=\;h(w_1),
\qquad
h_{i+1}\;=\;Y_{i+1}-i\,w_1^{\,i-1}Y_2+(i-1)\,w_1^{\,i}Y_1 ,
\]
so $(w_2,\dots,w_{n-2})=H^{-1}(w_1;h(w_1))$ are polynomial functions of $w_1$,
and the last remaining equation is the univariate \emph{fiber polynomial}
\[
R_Y(w_1)\;:=\;G_2\bigl(w_1,\,H^{-1}(w_1;h(w_1))\bigr)-\bigl(Y_2-w_1Y_1\bigr)
\;=\;0 .
\]
The fibers of the twisted map over targets with nonzero first coordinate are
in bijection with the roots of $R_Y$ at which $\gamma\neq0$. For $n=3$, $R_Y$
is the tangency polynomial of \S\ref{sec:sweep}.
\end{proposition}

\begin{proof}[Proof sketch]
The Jacobian and divisibility statements are verified exactly as for $F_6$
(the constancy was also confirmed at $8$ random rational points). For the
fibers, Proposition~\ref{prop:curvefib} gives the values of $w_2$ and $w_3$
along the fiber in closed form: writing them as $p_1(w_1)$ and $p_2(w_1)$,
\[
p_2(w_1)=Y_4-3Y_2w_1^2+2Y_1w_1^3-w_1^4,\qquad
p_1(w_1)=Y_3-2Y_2w_1+Y_1w_1^2-w_1^3-w_1^2\,p_2(w_1),
\]
\[
R_Y(w_1)\;=\;p_1(w_1)^2+p_1(w_1)+w_1\,p_2(w_1)+w_1^3+Y_1w_1-Y_2 ,
\]
of degree exactly $12$ with constant leading coefficient $1$ (the identity
$R_Y=$ residual of the second sweep equation was checked at $150$ random
$(Y,w_1)$). At the rational witness target $Y=(1,2,-1,3)$ the polynomial
$R_Y$ is squarefree and coprime to $\gamma(w_1)=Y_1-X_1(w_1,p_1(w_1),p_2(w_1))$
--- both certified by exact univariate gcd computations --- so that fiber has
exactly $12$ points, and by openness of both conditions the generic count is
$12$. Over $Y=0$,
\[
R_0(w_1)\;=\;w_1^{\,5}\,(w_1-1)\,
\bigl(w_1^6+w_1^5+w_1^4-w_1^3-w_1^2-w_1+1\bigr),
\]
the quintuple root $w_1=0$ has $\gamma=0$ and is deleted by the twist, while
$\gamma(1)=1$ and $\gcd(\gamma,R_0/w_1^5)=1$ exactly, so all seven remaining
roots survive: the fiber over $(C_0,0,0,0,0)$ has exactly $7$ points,
including the visible fixed point $(C_0,0,0,0,0)$.
\end{proof}

\end{document}
